\chapter{Fine-tuning LLMs}
\label{ch:finetuning}

\begin{quote}
\emph{``Pre-training gives the model knowledge. Fine-tuning gives it a job.''}
\end{quote}

% ============================================================
\section{Why fine-tune?}

Prompt engineering has limits. When you need a model to:
\begin{itemize}
  \item Consistently follow a specific output format,
  \item Handle domain-specific terminology (legal, medical, scientific),
  \item Achieve maximum accuracy on a narrow task,
  \item Reduce latency by using a smaller, specialized model,
\end{itemize}
fine-tuning is the answer.

% ============================================================
\section{Full fine-tuning vs parameter-efficient methods}

\textbf{Full fine-tuning} updates all model parameters. For a 7B model, this requires:
\begin{itemize}
  \item 28 GB for model weights (FP32)
  \item 28 GB for optimizer states (Adam)
  \item 28 GB for gradients
  \item Total: $\sim$84 GB VRAM --- multiple A100 GPUs.
\end{itemize}

\textbf{Parameter-efficient fine-tuning (PEFT)} freezes most parameters and trains a small number of additional ones.

% ============================================================
\section{LoRA: Low-Rank Adaptation}

LoRA decomposes weight updates into low-rank matrices. Instead of updating $W \in \mathbb{R}^{d \times d}$, we learn $\Delta W = BA$ where $B \in \mathbb{R}^{d \times r}$ and $A \in \mathbb{R}^{r \times d}$, with rank $r \ll d$.

\[
W' = W + \alpha \cdot BA
\]

Benefits: only $2dr$ parameters instead of $d^2$. For $d=4096, r=16$: from 16.7M to 131K parameters per layer.

\begin{minted}{python}
from peft import LoraConfig, get_peft_model, TaskType
from transformers import AutoModelForCausalLM

model = AutoModelForCausalLM.from_pretrained(
    "TinyLlama/TinyLlama-1.1B-Chat-v1.0",
    torch_dtype="auto",
    device_map="auto",
)

lora_config = LoraConfig(
    task_type=TaskType.CAUSAL_LM,
    r=16,                        # rank
    lora_alpha=32,               # scaling factor
    lora_dropout=0.05,
    target_modules=["q_proj", "v_proj", "k_proj", "o_proj"],
)

peft_model = get_peft_model(model, lora_config)
peft_model.print_trainable_parameters()
# trainable params: ~4M / total: ~1.1B = 0.36%
\end{minted}

% ============================================================
\section{QLoRA: quantized LoRA}

QLoRA loads the base model in 4-bit precision (NF4 quantization), then trains LoRA adapters in FP16. This reduces VRAM usage by $\sim$4x:

\begin{minted}{python}
from transformers import BitsAndBytesConfig
import torch

bnb_config = BitsAndBytesConfig(
    load_in_4bit=True,
    bnb_4bit_quant_type="nf4",
    bnb_4bit_compute_dtype=torch.float16,
    bnb_4bit_use_double_quant=True,
)

model = AutoModelForCausalLM.from_pretrained(
    "TinyLlama/TinyLlama-1.1B-Chat-v1.0",
    quantization_config=bnb_config,
    device_map="auto",
)

# Now apply LoRA on top of the quantized model
peft_model = get_peft_model(model, lora_config)
peft_model.print_trainable_parameters()
\end{minted}

\begin{aitip}
QLoRA makes fine-tuning a 7B model possible on a single free Colab T4 GPU (16 GB VRAM). This democratizes fine-tuning --- you do not need expensive hardware.
\end{aitip}

% ============================================================
\section{Datasets for fine-tuning}

\begin{longtable}{p{3cm}p{3cm}p{6.5cm}}
\toprule
\textbf{Dataset} & \textbf{Size} & \textbf{Description} \\
\midrule
Alpaca & 52K & Instruction-following data from Stanford \\
Dolly & 15K & Databricks employee-written instructions \\
OpenAssistant & 160K & Multi-turn conversation data \\
UltraChat & 1.5M & Synthetic multi-turn dialogues \\
\bottomrule
\end{longtable}

\begin{minted}{python}
from datasets import load_dataset

# Load Dolly dataset
dataset = load_dataset("databricks/databricks-dolly-15k", split="train")
print(f"Dataset size: {len(dataset)}")
print(dataset[0])

# Format for instruction tuning
def format_instruction(example):
    if example["context"]:
        text = (f"### Instruction:\n{example['instruction']}\n\n"
                f"### Context:\n{example['context']}\n\n"
                f"### Response:\n{example['response']}")
    else:
        text = (f"### Instruction:\n{example['instruction']}\n\n"
                f"### Response:\n{example['response']}")
    return {"text": text}

formatted = dataset.map(format_instruction)
print(formatted[0]["text"][:300])
\end{minted}

% ============================================================
\section{Fine-tuning with HuggingFace Trainer}

\begin{minted}{python}
from transformers import (
    AutoModelForCausalLM, AutoTokenizer,
    TrainingArguments, Trainer, DataCollatorForLanguageModeling,
    BitsAndBytesConfig,
)
from peft import LoraConfig, get_peft_model, TaskType
from datasets import load_dataset
import torch

# 1. Load model in 4-bit
bnb_config = BitsAndBytesConfig(
    load_in_4bit=True,
    bnb_4bit_quant_type="nf4",
    bnb_4bit_compute_dtype=torch.float16,
)

model_name = "TinyLlama/TinyLlama-1.1B-Chat-v1.0"
tokenizer = AutoTokenizer.from_pretrained(model_name)
tokenizer.pad_token = tokenizer.eos_token

model = AutoModelForCausalLM.from_pretrained(
    model_name, quantization_config=bnb_config, device_map="auto"
)

# 2. Apply LoRA
lora_config = LoraConfig(
    task_type=TaskType.CAUSAL_LM, r=16, lora_alpha=32,
    lora_dropout=0.05,
    target_modules=["q_proj", "v_proj", "k_proj", "o_proj"],
)
model = get_peft_model(model, lora_config)

# 3. Prepare dataset
dataset = load_dataset("databricks/databricks-dolly-15k", split="train[:1000]")

def tokenize(example):
    text = f"### Instruction:\n{example['instruction']}\n### Response:\n{example['response']}"
    return tokenizer(text, truncation=True, max_length=512, padding="max_length")

tokenized = dataset.map(tokenize, remove_columns=dataset.column_names)

# 4. Training arguments
training_args = TrainingArguments(
    output_dir="./tinyllama-dolly-lora",
    num_train_epochs=3,
    per_device_train_batch_size=4,
    gradient_accumulation_steps=4,
    learning_rate=2e-4,
    fp16=True,
    logging_steps=10,
    save_strategy="epoch",
    warmup_ratio=0.03,
    lr_scheduler_type="cosine",
)

# 5. Train
trainer = Trainer(
    model=model,
    args=training_args,
    train_dataset=tokenized,
    data_collator=DataCollatorForLanguageModeling(tokenizer, mlm=False),
)
trainer.train()

# 6. Save LoRA adapter
model.save_pretrained("./tinyllama-dolly-lora")
\end{minted}

% ============================================================
\section{Evaluation after fine-tuning}

\begin{minted}{python}
# Load the fine-tuned model
from peft import PeftModel

base_model = AutoModelForCausalLM.from_pretrained(
    model_name, quantization_config=bnb_config, device_map="auto"
)
finetuned = PeftModel.from_pretrained(base_model, "./tinyllama-dolly-lora")

# Compare base vs fine-tuned
prompt = "### Instruction:\nExplain what a neural network is.\n### Response:\n"
inputs = tokenizer(prompt, return_tensors="pt").to(model.device)

# Base model
base_out = base_model.generate(**inputs, max_new_tokens=100, temperature=0.7)
print("=== Base model ===")
print(tokenizer.decode(base_out[0], skip_special_tokens=True))

# Fine-tuned model
ft_out = finetuned.generate(**inputs, max_new_tokens=100, temperature=0.7)
print("\n=== Fine-tuned model ===")
print(tokenizer.decode(ft_out[0], skip_special_tokens=True))
\end{minted}

\begin{warningbox}
Fine-tuning on small datasets risks \emph{catastrophic forgetting}: the model loses general capabilities. Use a held-out validation set and monitor both task-specific and general performance.
\end{warningbox}

\begin{exercise}
\begin{enumerate}
  \item Fine-tune TinyLlama on 500 examples from the Dolly dataset using QLoRA. Report training loss at each epoch.
  \item Experiment with LoRA rank values: $r \in \{4, 8, 16, 32\}$. Which gives the best trade-off between quality and training time?
  \item Fine-tune on a custom dataset: create 100 instruction-response pairs about a topic you care about (e.g., cooking, history, or a specific domain). Evaluate the model before and after.
  \item Compare the VRAM usage of full fine-tuning vs QLoRA on TinyLlama. Use \texttt{torch.cuda.max\_memory\_allocated()}.
\end{enumerate}
\end{exercise}

\begin{ethicsbox}
Fine-tuning can embed biases from training data. If your dataset contains stereotypes, the model will learn them. Always audit your training data for bias, test the fine-tuned model on diverse inputs, and document the data sources and known limitations.
\end{ethicsbox}

% ============================================================
\section{Chapter summary}

\begin{itemize}
  \item Fine-tuning adapts a pre-trained model to a specific task or domain.
  \item Full fine-tuning updates all parameters and requires massive VRAM.
  \item LoRA trains low-rank adapter matrices, reducing trainable parameters by 99\%+.
  \item QLoRA combines 4-bit quantization with LoRA, enabling fine-tuning on free Colab GPUs.
  \item Instruction-tuning datasets (Alpaca, Dolly) teach models to follow instructions.
  \item HuggingFace Trainer + PEFT provides a complete fine-tuning pipeline in $\sim$30 lines.
  \item Always evaluate on held-out data and watch for catastrophic forgetting.
\end{itemize}
